\section{\framework}
\label{sec:framework}

\subsection{The task tuple}
A specification-diagnostic task is a tuple
\[
T = (M,\; R_p,\; J_t,\; C,\; M_{\mathrm{mon}},\; \Delta,\; \Pi),
\]
where $M$ is the environment dynamics, $R_p$ the proxy reward the agent is
trained on, $J_t$ the true objective it is judged by, $C$ the corruption channel
through which $R_p$ and $J_t$ differ, $M_{\mathrm{mon}}$ the information
available to a monitor, $\Delta$ a family of distribution shifts, and $\Pi$ the
policy class and training algorithm. Writing the tuple out forces each claim to
name its assumptions: a monitor result is a statement about a particular
$M_{\mathrm{mon}}$, a robustness result about a particular $\Delta$.

The design rule we impose on every task is that \emph{the corruption channel is
explicit and scalable}. In \PF{} and \ST{}, $R_p = J_t + \lambda\,C$ exactly; in
\GT{}, $R_p$ additionally omits the deadline that $J_t$ enforces. Scaling
$\lambda$ (the \emph{channel scale}) changes only the incentive and leaves the
dynamics untouched, so the channel's causal role can be tested by retraining
with a different $\lambda$ rather than inferred from correlations. Hack labels
are computed from behavior (how often the agent used the channel) and do not
depend on $\lambda$, so switching the channel off cannot hide the behavior from
the measurement (Section~\ref{sec:pitfalls}, item 3).

\subsection{Three tasks, three mechanisms}
All tasks are fully observed $7{\times}7$ gridworlds with walls and five actions
(four moves and \emph{interact}). Every episode samples a fresh layout, so an
agent cannot memorize a route, and every evaluation uses layouts the agent has
not been trained on. Observations are one-hot object planes plus a few scalars
(elapsed time and task state). Both rewards share a $-0.01$ step cost and pay
$+1$ for completing the intended task. Figure~\ref{fig:tasks} shows example
layouts and Table~\ref{tab:tasks} summarizes the channels.

\paragraph{\PF{} (farmable shaping).} The agent should reach a goal. The proxy
adds a bonus for every step spent on a beacon, a shaping term that is not
potential-based and therefore creates a reward cycle \citep{ng1999policy}.
Idling on the beacon earns more proxy reward than finishing the task. The OOD
split doubles wall density.

\paragraph{\ST{} (reward tampering).} The agent should pick up an object and
carry it to a target. The proxy is computed from a delivery sensor, and pressing
\emph{interact} on a console spoofs the sensor for four steps, a minimal
instance of an agent manipulating its reward process \citep{everitt2021reward}.
The OOD split delays the spoofing effect by two steps after each press.

\paragraph{\GT{} (omitted constraint).} The agent should reach a goal within a
deadline of the shortest-path length plus two steps. The proxy pays for
reaching the goal at any time and adds a bonus per gem. Gems are placed only
where the detour makes the deadline infeasible, so every gem collected trades
the true objective for proxy reward. The OOD split raises the number of gems
from three to five.

\begin{figure}[t]
\centering
\includegraphics[width=0.82\linewidth]{figures/tasks.pdf}
\caption{Example in-distribution (top) and OOD (bottom) layouts. A = agent,
G = goal or delivery target, dark cells are walls. Layouts are resampled every
episode.}
\label{fig:tasks}
\end{figure}

\begin{table}[t]
\centering
\caption{Corruption channels, shifts and hack labels. The full proxy reward and
true objective of each task are listed in Appendix~\ref{app:details}.}
\label{tab:tasks}
\small
\input{generated/table_tasks}
\end{table}
